\chapter{The Return and Risk of a Financial Portfolio}
\label{chap:The Return and Risk of a Financial Portfolio}

\begin{center}
\begin{tabular}{p{0.06\linewidth} | p{0.88\linewidth}}
	\toprule
	\multicolumn{2}{l}{\textbf{LEARNING OUTCOMES}} \\
	\midrule
	a & calculate, interpret, and evaluate the expected return, variance, standard deviation, covariance, and correlation of portfolio returns \\
	\midrule
	b & describe, calculate, and interpret the minimum-variance portfolio and portfolios that lie on the efficient frontier \\
	\midrule
	c & explain the selection of an optimal portfolio, given an investor's risk aversion and the capital allocation line, and how this extends to the market portfolio and the capital market line \\
	\bottomrule
\end{tabular}
\end{center}

\vspace{1em}

\begin{itemize}
	\item[] \textbf{Calculating Portfolio Statistics}
	      \begin{itemize}
			\item[] Returns of a Portfolio
			\item[] Variance and Standard Deviation of Portfolio Returns
			\item[] The Benefits of Diversification
			\item[] Types of Returns: Financial Indicators
			\item[] Long-Term Financial Asset Returns
			\item[] Risks and Returns
	      \end{itemize}
	\item[] \textbf{The Minimum-Variance Portfolio and the Efficient Frontier}
	      \begin{itemize}
			\item[] The Minimum-Variance Portfolio
			\item[] The Minimum-Variance Frontier
			\item[] Portfolio Optimization
	      \end{itemize}
	\item[] \textbf{Optimal Portfolio Selection}
	      \begin{itemize}
			\item[] The Capital Allocation Line (CAL)
			\item[] The Capital Market Line (CML)
			\item[] The Capital Asset Pricing Model (CAPM)
			\item[] Limitations
	      \end{itemize}
\end{itemize}
